\documentclass{article}
\usepackage[utf8]{inputenc}
\usepackage{amsmath}
\usepackage{geometry}
\geometry{margin=2cm}
\begin{document}

\section*{Questão 126 — ENEM 2024 — Ciências da Natureza}

\subsection*{Modos de transferência de calor em aquecedores solares}

A questão requer identificar qual processo físico transfere o calor do coletor solar para o reservatório térmico no aquecedor solar apresentado.

\subsection*{Interpretação e Estratégia}

Deve-se compreender os modos clássicos de transferência de calor: condução, convecção, radiação, difusão e absorção. Observando o esquema, fica evidente que a água circula do coletor para o reservatório, indicando transferência pelo fluido.

\subsection*{Análise dos Conteúdos e Mini Revisão}

\begin{description}
  \item[Condução:] Transferência de calor por contato direto, predominante em sólidos.
  \item[Convecção:] Transferência de calor através do movimento de fluidos, neste caso a água quente que circula.
  \item[Radiação:] Transferência de energia por ondas eletromagnéticas, responsável pelo aquecimento inicial do coletor pelo sol.
  \item[Difusão:] Movimentação de partículas que não caracteriza modo principal de transferência térmica nesse contexto.
  \item[Absorção:] Captação da energia solar pelo coletor, não transferência interna.
\end{description}

\subsection*{Resolução e "Pulo do Gato"}

A radiação solar incide no coletor de tubo a vácuo e aquece a água. Esta água aquecida sobe para o reservatório, enquanto a água fria desce para o coletor, estabelecendo uma circulação natural. Esse movimento da água quente é convecção, que transporta o calor para o reservatório.

\textbf{Pulo do gato:} Radiação aquece o coletor, mas o calor é transferido internamente por convecção, pelo movimento da água aquecida.

\subsection*{Armadilhas e Dificuldades Comuns}

Cuidado para não confundir difusão, absorção, condução e radiação com o modo físico predominante de transferência de calor do coletor para o reservatório.

\subsection*{Código da Questão}

CNC\_FIS\_TRANSFERENCIA\_001

\end{document}


\end{document}